\documentclass[11pt]{article}

\usepackage[margin=1in]{geometry}
\usepackage{enumitem}
\usepackage{listings}

\lstdefinestyle{shell}{
  basicstyle=\ttfamily\small,
  columns=fullflexible,
  keepspaces=true,
  frame=single,
  breaklines=true,
  showstringspaces=false
}

\title{Yalnix Checkpoint 4 Test Notes}
\author{COSC 58}
\date{\today}

\begin{document}
\maketitle

\section{Test Commands}

Run from the Thayer source directory:

\begin{lstlisting}[style=shell]
make clean
make
timeout 8 ./yalnix -W cp4_fork
timeout 8 ./yalnix -W cp4_exec
timeout 8 ./yalnix -W cp4_stack
\end{lstlisting}

\section{Expected Observations}

\begin{itemize}[noitemsep]
  \item \texttt{cp4\_fork} forks two children, each child exits with a different status, and the parent collects both with \texttt{Wait}.
  \item A third \texttt{Wait} in \texttt{cp4\_fork} returns \texttt{ERROR}, showing the no-remaining-children case.
  \item \texttt{cp4\_exec} replaces itself with \texttt{cp4\_target}; the pid remains the same and argc/argv are rebuilt by \texttt{LoadProgram}.
  \item \texttt{cp4\_stack} touches more stack than the initially mapped user stack page, causing \texttt{TRAP\_MEMORY} stack growth before exiting.
  \item Each test exits as the initial process, which halts the machine cleanly.
\end{itemize}

\section{Run Evidence}

Representative \texttt{-W} trace lines:

\begin{lstlisting}[style=shell]
cp4_fork_W_RC=0
User Prog cp4_fork: parent pid 1
Hardware  | Syscall trap Fork
Hardware  | Syscall trap Fork
Hardware  | Syscall trap Wait
User Prog cp4_fork: child A pid 2
User Prog cp4_fork: child B pid 3
User Prog cp4_fork: wait pid 3 status 22
User Prog cp4_fork: wait pid 2 status 11
User Prog cp4_fork: final wait rc -1

cp4_exec_W_RC=0
User Prog cp4_exec: before exec pid 1
Hardware  | Syscall trap Exec
User Prog cp4_target: pid 1 argc 3
User Prog cp4_target: argv cp4_target alpha beta

cp4_stack_W_RC=0
Hardware  | Calling Yalnix handler for trap TRAP_MEMORY
User Prog cp4_stack: touched 24576 bytes
\end{lstlisting}

\end{document}
